Los resultados obtenidos demuestran que se logró el objetivo general propuesto para esta investigación. Se desarrolló un sistema de información que integra los resultados individuales de las pruebas Saber 11 y Saber Pro de la Universidad de San Buenaventura, permitiendo estimar el desempeño esperado de los estudiantes mediante una red neuronal y calcular el valor agregado institucional en cada competencia y periodo analizado. Adicionalmente, la información generada por el modelo fue incorporada en un tablero interactivo que facilita el acceso a los resultados y proporciona una herramienta de apoyo para los procesos de análisis y toma de decisiones dentro de la institución.
No obstante, la hipótesis planteada no se confirma en ninguna de sus dos afirmaciones: la capacidad predictiva máxima alcanzada $(R^2 \approx 0,61)$ no llega al umbral del $70\%$ propuesto, y el análisis del valor agregado no evidencia una mejora promedio del $10\%$, sino un aporte institucional mayoritariamente negativo. Sin embargo, esto no compromete la utilidad del sistema, este cumple su función analítica con independencia de que los datos confirmen o no la expectativa inicial, dado que la función analítica del sistema no es arrojar un valor esperado de antemano, sino integrar la información, estimar el desempeño y cuantificar el valor agregado de forma válida y verificable. Desde esta perspectiva, puede afirmarse que el sistema alcanzó el propósito para el cual fue desarrollado. La herramienta permitió consolidar cerca de $1,32$ millones de pares de estudiantes, generar estimaciones validadas y con error cuantificado, medir el valor agregado por competencia y por año, y poner esta información a disposición de la institución mediante un entorno de consulta accesible. Los resultados relacionados con el valor agregado no deben interpretarse como una limitación del sistema. Por el contrario, constituyen uno de los hallazgos más relevantes del estudio, ya que permiten dimensionar una situación institucional que difícilmente habría podido identificarse con el mismo nivel de detalle sin una herramienta de este tipo. El hecho de reportar estos resultados de manera transparente, incluso cuando difieren de las expectativas iniciales, evidencia la utilidad del sistema como apoyo para el análisis y la toma de decisiones basadas en datos. 
Entre los modelos evaluados, la red neuronal mostró el mejor desempeño general, manteniendo una ventaja pequeña pero consistente frente a la regresión lineal y al algoritmo \textit{K-means} en la mayoría de las competencias. La única excepción fue el módulo de Inglés, donde los resultados sugieren que una relación de naturaleza más lineal permite alcanzar un desempeño similar o incluso ligeramente superior. Asimismo, se encontró que las variables socioeconómicas incluidas en el estudio no aportan una mejora significativa en la capacidad predictiva una vez que se conocen los resultados de Saber 11. Este hallazgo representa un aporte metodológico importante para futuras investigaciones sobre esta población y ofrece una referencia útil para el diseño de modelos predictivos similares.
Por último, el aporte institucional negativo, lejos de invalidar el proyecto, constituye un insumo accionable: señala competencias y periodos en los que el desempeño observado quedó por debajo del esperado y orienta la atención hacia las áreas que requieren intervención. El valor de este trabajo reside, por tanto, en la herramienta construida y en la evidencia que produce, más que en la verificación de la hipótesis.